% =========================================================
%  SECCIÓN V — DISCUSIÓN
% =========================================================
\section{Discusi\'on}

\subsection{Comparaci\'on entre Versiones de Vectorizaci\'on}

Los resultados muestran que la vectorizaci\'on, aunque presente,
produce ganancias modestas sobre el kernel Prewitt en el hardware
utilizado.
La versi\'on autom\'atica es la que mejor rendimiento ofrece
($1.24\times$), seguida de la expl\'icita ($1.21\times$).
Resulta llamativo que la vectorizaci\'on guiada con
\texttt{simdlen(8)} no mejore significativamente sobre la escalar
($1.01\times$) pese a estar vectorizada con AVX2.

Este comportamiento se explica por el papel del ancho vectorial
real.
La versi\'on autom\'atica utiliza un ancho de~16~floats (AVX2
con posible unroll adicional), procesando el doble de datos por
iteraci\'on que la guiada (ancho~8).
La expl\'icita, aunque nominalmente de ancho~8, recibe un
desenrollado $4\times$ del compilador, resultando en un ancho
efectivo de~32, lo que tambi\'en supera a la guiada.
Sin embargo, todos estos anchos vectoriales son incapaces de
saturar el ancho de banda de memoria en igual medida que lo
har\'ian en un kernel de mayor intensidad aritm\'etica.

Desde la perspectiva de la complejidad de implementaci\'on:
\begin{itemize}
  \item La \textbf{autom\'atica} requiere el menor esfuerzo y
    produce el mejor resultado; solo necesita compilar con
    \texttt{-O3 -march=native} y usar \texttt{\_\_restrict}.
  \item La \textbf{guiada} requiere conocer el pragma de OpenMP
    SIMD y elegir el \texttt{simdlen} adecuado; result\'o inferior
    porque \texttt{simdlen(8)} es m\'as conservador que la
    decisi\'on autom\'atica del compilador.
  \item La \textbf{expl\'icita} requiere dominar los intr\'insecos
    AVX2 y resulta en rendimiento ligeramente menor que la
    autom\'atica, mostrando que el compilador tambi\'en puede
    optimizar el c\'odigo escalar mejor que la versi\'on manual
    de ancho~8 sin unroll.
\end{itemize}

\subsection{Interpretaci\'on del Modelo Roofline}

La intensidad aritm\'etica te\'orica del kernel es
$\mathcal{I} = 0.333$ FLOP/Byte (ec.~\ref{eq:ai}).
Este valor sit\'ua el kernel \textit{por debajo} del punto de
cresta del Roofline calculado como:

\begin{equation}
  \mathcal{I}_{\text{cresta}} =
  \frac{P_{\text{peak}}}{B_{\text{mem}}} \approx
  \frac{9.28 \text{ GFLOPS}}{24 \text{ GB/s}} \approx
  0.387 \text{ FLOP/Byte}
\end{equation}

Al ser $\mathcal{I} = 0.333 < \mathcal{I}_{\text{cresta}} = 0.387$,
el kernel es \textbf{memory-bound}: el cuello de botella es el
ancho de banda de memoria, no la capacidad de c\'omputo.
Esto significa que aunque se utilicen instrucciones AVX2 de
m\'axima eficiencia, los registros YMM permanecer\'an ociosos
una parte del tiempo esperando que la memoria entregue los datos.

El techo alcanzable con la memoria DRAM a $24$~GB/s ser\'ia:
\begin{equation}
  P_{\max}^{\text{DRAM}} = 0.333 \times 24 \approx 8.0 \text{ GFLOPS}
\end{equation}

Los resultados experimentales (5.92--7.34~GFLOPS) son coherentes
con este l\'imite te\'orico, confirmando que el kernel opera
pr\'oxima a la banda de datos de la jerarqu\'ia de cach\'e
L2/L3 del procesador ($65$--$190$~GB/s) m\'as que de la DRAM,
lo que explica que los GFLOPS medidos superen el techo DRAM.

\subsection{ISA Detectada por Versi\'on}

Intel Advisor reporta la ISA (\textit{Instruction Set Architecture})
utilizada en cada programa:

\begin{itemize}
  \item \textbf{Escalar}: AVX2, AVX, SSE2, \textbf{SSE}.
    La presencia de instrucciones SSE puras (sin empaquetado)
    confirma que el kernel usa registros XMM en modo escalar
    (un solo float por registro), no SIMD empaquetado.
    El flag \texttt{-no-vec} proh\'ibe el uso vectorial, pero
    el procesador sigue usando los registros SSE para operaciones
    escalares.

  \item \textbf{Autom\'atica}: AVX2, AVX, SSE2 (sin SSE puro).
    La ausencia de SSE puro confirma que el kernel escal\'o al
    modo AVX2 empaquetado.

  \item \textbf{Guiada}: AVX2, AVX, SSE2 (sin SSE puro).
    Idéntico al autom\'atico en cuanto a ISA, aunque con ancho~8
    en vez de~16.

  \item \textbf{Expl\'icita}: AVX2, AVX, SSE2 (sin SSE puro).
    Los intr\'insecos \texttt{\_mm256\_*} garantizan el uso de
    registros YMM de 256~bits.
\end{itemize}

\subsection{Dominancia del C\'odigo Escalar y Ley de Amdahl}

Un resultado relevante es que el bucle de replicaci\'on de
datos (el que extiende la imagen 800 veces en memoria para
alcanzar $N_{\text{tot}}$) consume entre $0.040$ y $0.070$~s,
superando al propio kernel vectorizado ($0.030$~s).
Esto implica que la mayor parte del tiempo lo ocupan operaciones
fuera del kernel: lectura del archivo PGM, c\'alculo de los
gradientes Prewitt con las m\'ascaras $3\times3$ (bucle escalar
de 254 iteraciones), replicaci\'on de datos y escritura
del resultado.

Este comportamiento puede analizarse formalmente con la
\textbf{Ley de Amdahl}~\cite{sutter2005}.
Sea $p$ la fracci\'on paralelizable del programa y $f$ el
factor de aceleraci\'on (\textit{speedup}) de esa fracci\'on.
El speedup total del programa est\'a acotado por:

\begin{equation}
  S(f) = \frac{1}{(1-p) + p/f}
  \label{eq:amdahl}
\end{equation}

En nuestro caso, $p \approx 0.12$ (el kernel representa el
$12\%$ del tiempo total) y $1-p \approx 0.88$ es la fracci\'on
escalar fija.
La Tabla~\ref{tab:amdahl} muestra el speedup total m\'aximo
te\'orico en funci\'on del factor $f$ del kernel:

\begin{table}[htb]
\centering
\caption{Speedup m\'aximo total predicho por la Ley de Amdahl
($p=0.12$, $1-p=0.88$).}
\label{tab:amdahl}
\footnotesize
\begin{tabular}{@{}ccc@{}}
\toprule
\textbf{Speedup kernel} $f$ & \textbf{Speedup total} $S(f)$ &
\textbf{Mejora (\%)} \\
\midrule
$1\times$ (escalar)  & 1.00 &  ---  \\
$2\times$            & 1.09 &  +9\% \\
$4\times$            & 1.11 & +11\% \\
$8\times$  (AVX2 t.) & 1.12 & +12\% \\
$16\times$ (AVX2 m.) & 1.12 & +12\% \\
$\infty$             & 1.14 & +14\% \\
\bottomrule
\end{tabular}
\end{table}

El speedup te\'orico m\'aximo posible, incluso vectorizando el
kernel infinitamente r\'apido, es de apenas $1.14\times$.
Los speedups experimentales observados ($1.01$--$1.24\times$)
son consistentes con este l\'imite te\'orico, validando que
el dise\~no del experimento es correcto.
La \'unica forma de superar este techo ser\'ia vectorizar
tambi\'en las partes escalares del programa (c\'alculo de
gradientes, replicaci\'on de datos).

\subsection{Limitaciones y Trabajo Futuro}

La m\'etrica \textit{Vectorization Gain/Efficiency} de Intel
Advisor reporta \textit{Not Available} en todas las versiones.
Esto se debe a que dicha m\'etrica requiere el compilador Intel
C++ cl\'asico (\texttt{icc}/\texttt{icpc} versi\'on 16+), y el
presente trabajo utiliz\'o \texttt{icpx} (compilador DPC++ de
la generaci\'on oneAPI).
Para trabajo futuro se recomienda explorar:
(a)~operadores de mayor intensidad aritm\'etica como el filtro
Gaussiano o la convolución 2D densa, que ser\'ian m\'as
sensibles a la vectorizaci\'on;
(b)~an\'alisis multi-hilo con OpenMP para aprovechar los
14~n\'ucleos disponibles;
(c)~uso de AVX-512 si el hardware lo soporta, para alcanzar
anchos de~16 o~32 floats con mayor eficiencia que el unroll.
